\documentclass[../../main.tex]{subfiles}
\begin{document}

\begin{flushright}
  \footnotesize\textit{【自動文字起こし・要確認】}
\end{flushright}

{\bf [解]}
点$X$に対し，$\overrightarrow{AX}=\vec x$とおく．$s,t,u$を$0<s,t,u<1$を満たす実数として
\begin{align}
  \vec p=s\vec b，\qquad\vec q=t\vec c+(1-t)\vec b，\qquad\vec r=u\vec c
\end{align}
とおける（$\vec b=\overrightarrow{AB}$，$\vec c=\overrightarrow{AC}$）．$\triangle PQR$の重心を$G$とする．
\begin{align}
  3\vec g&=(t+u)\vec c+(1+s-t)\vec b
\end{align}

ここで，$\lambda=1+s-t$，$\mu=t+u$とおくと，$s,t,u\in(0,1)$が独立に動くことから，$(\lambda,\mu)$は頂点
\begin{align}
  (1,0),(2,0),(2,1),(1,2),(0,2),(0,1)
\end{align}
を持つ六角形の内部（境界含まず）を動く．$\vec g=\dfrac{\lambda}{3}\vec b+\dfrac{\mu}{3}\vec c$だから，これは，各辺を$1:2$，$2:1$に内分する6点を頂点とする六角形の内部に対応する．

だから，下図着色部（境界含まず）．

\begin{figure}[htb]
  \centering
  \begin{tikzpicture}[x=2cm, y=2cm, >=stealth,
    every node/.style={font=\small}]
    \begin{scope}
      \coordinate (A) at (1,1.7);
      \coordinate (B) at (0,0);
      \coordinate (C) at (1.6,0);
      \draw[line width=1.25pt, blue!85!black] (A) -- (B) -- (C) -- cycle;
      \node[above] at (A) {$A$};
      \node[below left] at (B) {$B$};
      \node[below right] at (C) {$C$};
    \end{scope}
    \begin{scope}[shift={(2.25,0)}]
      \coordinate (A) at (1,1.7);
      \coordinate (B) at (0,0);
      \coordinate (C) at (1.6,0);
      \coordinate (P1) at ($(A)!1/3!(B)$);
      \coordinate (P2) at ($(A)!2/3!(B)$);
      \coordinate (Q1) at ($(B)!1/3!(C)$);
      \coordinate (Q2) at ($(B)!2/3!(C)$);
      \coordinate (R1) at ($(C)!1/3!(A)$);
      \coordinate (R2) at ($(C)!2/3!(A)$);
      \fill[blue!15] (P1) -- (P2) -- (Q1) -- (Q2) -- (R1) -- (R2) -- cycle;
      \draw[line width=1.25pt, blue!85!black] (A) -- (B) -- (C) -- cycle;
      \draw[dashed, line width=1pt, teal!80!black]
      (P1) -- (P2) -- (Q1) -- (Q2) -- (R1) -- (R2) -- cycle;
      \foreach \p in {P1,P2,Q1,Q2,R1,R2}
      \draw[teal!80!black, fill=white] (\p) circle (1.3pt);
      \node[above] at (A) {$A$};
      \node[below left] at (B) {$B$};
      \node[below right] at (C) {$C$};
    \end{scope}
  \end{tikzpicture}
  \caption{三角形$ABC$と重心$G$が動く範囲（着色部，境界含まず）}
\end{figure}

\newpage

\end{document}
